\section*{Chapitre \ref{ch:nw:programmable-hardware-i-architecture-and-toolchain} --- Matériel programmable I~: architecture et chaîne d'outils}

\begin{solution}{exo:nw:programmable-hardware-i-architecture-and-toolchain:1}
$25\,000/32 = 781.25$~MHz à 32 bits~; $25\,000/128 = 195.3125$~MHz à 128 bits.
\end{solution}

\begin{solution}{exo:nw:programmable-hardware-i-architecture-and-toolchain:2}
Période $12 \times 0.5 + 0.4 = \qty{6.4}{\nano\second}$~: 156.25~MHz, et un cycle de latence, soit \qty{6.4}{\nano\second}.
\end{solution}

\begin{solution}{exo:nw:programmable-hardware-i-architecture-and-toolchain:3}
Les octets 12 à 15 se trouvent dans le battement~1, comme les derniers octets 6 à 9~: on reste au cycle~2. Les octets 14 à 17 se terminent dans le battement~2 (octets 16 à 23), accepté au
cycle~2~: le champ est enregistré au cycle~3.
\end{solution}

\begin{solution}{exo:nw:programmable-hardware-i-architecture-and-toolchain:4}
La boucle SystemVerilog est déroulée en huit voies d'octets qui existent côte à côte~; chaque voie est évaluée à chaque cycle, et le résultat est
sélectionné par des multiplexeurs dont le délai est fixe. La boucle C++ exécute ses itérations les unes après les autres, avec des branchements dont le coût dépend des
données et du prédicteur, et sur une machine partagée avec d'autres travaux.
\end{solution}

\begin{solution}{exo:nw:programmable-hardware-i-architecture-and-toolchain:5}
Six étages laissent $\lceil 12/6\rceil = 2$ niveaux par étage, cinq en laissent 3, comme quatre~: seul l'étage le plus profond compte, si bien que les niveaux doivent
se répartir également entre les étages pour qu'un nouvel étage apporte un gain. Les concepteurs comptent les niveaux sur le \omterm{def:nw:programmable-hardware-i-architecture-and-toolchain:timing}{chemin critique}, et non en totalité.
\end{solution}

\begin{solution}{exo:nw:programmable-hardware-i-architecture-and-toolchain:6}
Il retient le battement déjà en vol lorsque l'aval se bloque, si bien que l'amont peut voir un signal «~prêt~» enregistré (aucun chemin combinatoire
à travers l'étage) sans perdre de données. Avec un seul registre et un signal «~prêt~» enregistré, le battement envoyé au cycle où le blocage est
constaté serait écrasé ou perdu.
\end{solution}

\begin{solution}{exo:nw:programmable-hardware-i-architecture-and-toolchain:7}
Les deux simulateurs concordent avec le modèle de référence. Le champ (octets 10 à 17) se termine dans le battement~2~; il est enregistré au cycle qui suit
celui où le tampon à saut enregistre ce battement, comme pour le champ du chapitre.
\end{solution}

\begin{solution}{exo:nw:programmable-hardware-i-architecture-and-toolchain:8}
Les deux simulateurs interprètent le langage de façon indépendante~; une conception qui repose sur un ordre que le langage laisse ouvert (affectations
bloquantes entre processus, lecture d'un signal au cycle même où il est écrit) peut réussir sur l'un et échouer sur l'autre, et échouer encore sur
la puce. Un désaccord est peu coûteux à trouver en simulation et coûteux sur la carte~; aucun des deux simulateurs ne remplace non plus l'analyse temporelle.
\end{solution}

\begin{solution}{pb:nw:programmable-hardware-i-architecture-and-toolchain:1}
\begin{enumerate}
\item 156.25~MHz et 390.625~MHz.
\item $1/6.4\,\mathrm{ns} = 156.25$~MHz.
\item La ligne à \qty{10}{\giga\bit\per\second} exactement, sans marge~; pas celle à \qty{25}{\giga\bit\per\second}.
\item 16 battements~: $16 \times 6.4 = \qty{102.4}{\nano\second}$ à 156.25~MHz, $16 \times 2.56 = \qty{40.96}{\nano\second}$ à 390.625~MHz.
\item Trois~: $\lceil 12/3\rceil = 4$ niveaux, une période de \qty{2.4}{\nano\second}.
\item 416.7~MHz.
\item $3 \times 2.56 = \qty{7.68}{\nano\second}$.
\item $3 \times 2.4 = \qty{7.2}{\nano\second}$.
\item L'octet 40 se trouve dans le sixième battement (octets 40 à 47)~: six battements.
\item Six battements et trois étages, neuf cycles à \qty{2.56}{\nano\second}~: \qty{23.04}{\nano\second}.
\item Neuf des seize battements sont arrivés~: sept battements, soit 56 octets, restent à venir.
\item La décision est prête \qty{18}{\nano\second} avant la fin de la trame~; attendre la séquence de contrôle ajouterait ces nanosecondes. Le risque est
  d'agir sur une trame corrompue, que la conception doit pouvoir annuler ou accepter comme rare.
\item \textbf{Résultat nommé.} Trois étages atteignent 390.625~MHz, et la décision prend \qty{7.68}{\nano\second} à l'horloge de la ligne, contre
  \qty{6.4}{\nano\second} pour la conception non pipelinée à sa propre fréquence de 156.25~MHz, qui ne peut suivre la ligne plus rapide à cette largeur.
\item Oui~: à 128 bits, la ligne exige 195.3125~MHz, que deux étages atteignent (294~MHz), pour une latence de $2 \times 5.12 = \qty{10.24}{\nano\second}$ à
  l'horloge de la ligne, en moins de battements, plus larges.
\item Le délai par niveau et par registre, et les délais de routage, chemin par chemin.
\item Chaque étage dure au moins autant que sa logique, et les périodes des étages sont fixées par le plus profond.
\item Parce qu'une conception qui ne peut suivre la ligne perd des données ou exige un chemin de données plus large, et que quelques nanosecondes de décision valent
  moins que de voir chaque message.
\item Exécuter les deux sur le même stimulus dans les deux simulateurs et comparer leurs sorties, décalées des cycles supplémentaires.
\item Un synchroniseur ou une FIFO entre les horloges~: quelques cycles de latence et une source d'erreurs rares, difficiles à reproduire si elle est mal faite.
\item Elle fixe l'horloge, et l'horloge fixe la quantité de logique que chaque cycle peut contenir.
\end{enumerate}
\end{solution}

\begin{solution}{iq:nw:programmable-hardware-i-architecture-and-toolchain:1}
Une puce de logique programmable, de mémoires et d'émetteurs-récepteurs configurée par un fichier~: la conception est un circuit qui traite le flux réseau
à mesure qu'il arrive, en un nombre fixe de cycles, sans système d'exploitation ni ordonnanceur sur le chemin. Les sociétés y recourent là où le déterminisme et les dernières
centaines de nanosecondes comptent~: traitement des flux de marché, filtrage, contrôles de risque et ordres déclenchés.
\iqlookfor{le déterminisme et le traitement au fil de l'eau, et non seulement la vitesse.}
\end{solution}

\begin{solution}{iq:nw:programmable-hardware-i-architecture-and-toolchain:2}
Faire tenir chaque chemin entre registres dans la période d'horloge. Sur un chemin en échec~: le pipeliner (ajouter un registre), restructurer la logique
(moins de niveaux, utiliser la chaîne de retenue ou une mémoire), réduire le fan-out, contraindre le placement, ou abaisser l'horloge si la ligne le permet.
\iqlookfor{le \omterm{def:nw:programmable-hardware-i-architecture-and-toolchain:timing}{chemin critique} et les remèdes concrets.}
\end{solution}

\begin{solution}{iq:nw:programmable-hardware-i-architecture-and-toolchain:3}
Il raccourcit chaque étage, ce qui élève l'horloge, de sorte que la conception suit une ligne plus rapide ou accepte une nouvelle entrée à chaque cycle. La latence
d'une décision ne baisse pas~; le débit et la capacité à tenir l'horloge de la ligne augmentent.
\iqlookfor{la latence face au débit et à l'horloge de la ligne.}
\end{solution}

\begin{solution}{iq:nw:programmable-hardware-i-architecture-and-toolchain:4}
Le RTL donne un contrôle au cycle près sur le \omterm{def:nw:programmable-hardware-i-architecture-and-toolchain:timing}{chemin critique} et sur la latence, dont les conceptions les plus rapides ont besoin. Le HLS s'écrit et se modifie plus vite, ouvre
la contribution des ingénieurs logiciels et permet d'explorer rapidement les compromis, au prix d'un contrôle moindre sur l'ordonnancement et d'un débogage plus difficile de ce que
le compilateur a produit. Beaucoup d'équipes utilisent le HLS pour les blocs les moins critiques et le RTL pour le \omterm{def:nw:programmable-hardware-i-architecture-and-toolchain:timing}{chemin critique} (hot path).
\iqlookfor{le contrôle face à la productivité, bloc par bloc.}
\end{solution}

\begin{solution}{iq:nw:programmable-hardware-i-architecture-and-toolchain:5}
Un modèle de référence logiciel~; simulation dans deux simulateurs sur des flux enregistrés et synthétiques, avec trous, rafales et trames mal formées,
avec une contre-pression (back-pressure) aléatoire~; vérifications de propriétés (aucun message perdu ou dupliqué)~; puis rejeu matériel en boucle des captures enregistrées et
comparaison de la sortie de la carte avec celle du gestionnaire logiciel.
\iqlookfor{un modèle de référence et le rejeu de captures réelles, pas seulement des tests unitaires.}
\end{solution}

\begin{solution}{iq:nw:programmable-hardware-i-architecture-and-toolchain:6}
Un signal qui passe entre des horloges sans lien~: échantillonné pendant qu'il change, une bascule peut entrer en métastabilité, et des bits d'un bus peuvent être
échantillonnés de part et d'autre d'un changement. Il faut un synchroniseur pour les bits isolés et une FIFO asynchrone ou une poignée de main pour les données~; les défaillances
sont rares et impossibles à reproduire.
\iqlookfor{la métastabilité et les structures usuelles.}
\end{solution}
